\section{Cronograma de Trabajo}

\subsection{Fase 1: Configuración y Extracción de Features (Sem.\ 1--3)}

\textbf{Objetivo:} Implementar un pipeline para la extracción de características
acústicas.

\textbf{Semana 1 — Configuración del entorno técnico:}
\begin{itemize}
    \item Instalar Python 3.8+, Librosa, OpenSMILE, pyannote.audio.
    \item Descargar datasets RAVDESS.
    \item Crear la estructura del proyecto y el repositorio de Git.
    \item Probar la extracción básica de audio con Librosa.
\end{itemize}

\textbf{Semana 2 — Implementación de extracción de features:}
\begin{itemize}
    \item Desarrollar una función para extraer la F0 con Librosa.
    \item Implementar cálculo de intensidad.
    \item Crear módulo de segmentación temporal.
\end{itemize}

\textbf{Semana 3 — Features de calidad vocal y espectrales:}
\begin{itemize}
    \item Implementar el cálculo de Jitter y Shimmer con Parselmouth.
    \item Extraer 13 MFCC por segmento.
    \item Calcular la tasa de habla.
    \item Desarrollar módulo de visualización.
\end{itemize}

\subsection{Fase 2: Baseline Emocional y Preprocesamiento (Sem.\ 4--5)}

\textbf{Objetivo:} Evaluar métodos para establecer líneas de base emocionales.

\textbf{Semana 4 — Implementación de métodos de baseline:}
\begin{itemize}
    \item Método 1: Baseline de inicio (primeros 2--5 minutos).
    \item Método 2: Baseline por mediana global.
    \item Método 3: Baseline adaptativo con ventana deslizante.
    \item Implementar la normalización de las features con respecto a la baseline.
\end{itemize}

\textbf{Semana 5 — Preprocesamiento avanzado de audio:}
\begin{itemize}
    \item Implementar diarización con pyannote.audio (separar asesor y cliente).
    \item Desarrollar filtros de ruido y de normalización de volumen.
    \item Crear un detector de silencios.
    \item Aplicar el preprocesamiento a un subconjunto del conjunto de datos y
          validar su calidad.
\end{itemize}

\subsection{Fase 3: Detección de Picos Emocionales (Sem.\ 6--8)}

\textbf{Objetivo:} Desarrollar algoritmos de detección de anomalías.

\textbf{Semana 6 — Implementación de detector estadístico:}
\begin{itemize}
    \item Calcular z-scores de features respecto a baseline.
    \item Implementar detección por umbral.
    \item Desarrollar un algoritmo de ventanas deslizantes para la detección temporal.
\end{itemize}

\textbf{Semana 7 — Implementación de detectores de ML:}
\begin{itemize}
    \item Entrenar Isolation Forest para la detección de anomalías multivariadas.
    \item Implementar One-Class SVM como método alternativo.
    \item Realizar un grid search para optimizar los hiperparámetros.
    \item Comparar el rendimiento del método estadístico vs.\ ML.
\end{itemize}

\textbf{Semana 8 — Integración y caracterización de picos:}
\begin{itemize}
    \item Crear un pipeline integrado que combine el baseline y la detección.
    \item Implementar sistema de etiquetado temporal.
\end{itemize}

\subsection{Fase 4: Clasificación de Valencia Emocional (Sem.\ 9--10)}

\textbf{Objetivo:} Clasificar los picos según su valencia emocional.

\textbf{Semana 9 — Desarrollo de clasificador de valencia:}
\begin{itemize}
    \item Preparar datos etiquetados de RAVDESS/IEMOCAP (positivos/negativos).
    \item Entrenar un clasificador binario SVM con características de picos detectadas.
    \item Implementar Random Forest como método alternativo.
\end{itemize}

\textbf{Semana 10 — Optimización:}
\begin{itemize}
    \item Optimizar hiperparámetros mediante grid search.
    \item Analizar la importancia de las features en la clasificación de la valencia.
    \item Probar el clasificador en llamadas reales de un call center.
\end{itemize}

\subsection{Fase 5: Validación y Experimentación (Sem.\ 11--13)}

\textbf{Objetivo:} Validar el sistema mediante llamadas reales y anotaciones de
expertos.

\textbf{Semana 11 — Validación:}
\begin{itemize}
    \item Ejecutar el sistema completo sobre las llamadas anotadas.
    \item Calcular métricas: Precision, Recall, F1-score, Accuracy.
    \item Crear la matriz de confusión y analizar las curvas ROC.
    \item Documentar las discrepancias entre las anotaciones y las detecciones
          automáticas.
\end{itemize}

\textbf{Semana 12 — Experimentación comparativa:}
\begin{itemize}
    \item Comparar el rendimiento de 3 métodos de baseline.
    \item Comparar detectores y clasificadores de valencia.
    \item Generar tablas y gráficos comparativos de los resultados.
\end{itemize}

\textbf{Semana 13 — Análisis de resultados:}
\begin{itemize}
    \item Analizar casos de éxito, falsos positivos y falsos negativos.
    \item Ajustar umbrales y parámetros con base en el análisis.
\end{itemize}

\subsection{Fase 6: Interfaz y Documentación Final (Sem.\ 14--15)}

\textbf{Objetivo:} Integrar todos los objetivos en el producto final.

\textbf{Semana 14 — Finalización y documentación:}
\begin{itemize}
    \item Optimizar y limpiar el código final.
    \item Documentar el código con docstrings y comentarios.
    \item Consolidar el documento final en el formato institucional.
    \item Verificar referencias IEEE.
    \item Preparar la presentación y la demo en vivo.
\end{itemize}
